% q0062.tex — The More, the Merrier?
\question{
  id        = {0062},
  title     = {The More, the Merrier?},
  source    = {OBECON},
  year      = {2026},
  round     = {Seletiva IEO -- Phase C},
  language  = {English},
  topic     = {Macro},
  subtopic  = {Real vs Nominal Income},
  type      = {objective},
  statement = {%
    An individual has real income $R_0$ at $t = 0$. Nominal income grows at
    $s_t\% = t + 5$ per year, while inflation is $\pi_t = t^2 - 25t + 155$.
    At which value of $t$ is real income $R_t$ maximised?
  },
  choices   = {%
    \item 8.
    \item 13.
    \item 17.
    \item None of the above.
  },
  answer    = {D},
  solution  = {%
    Real income grows when nominal growth exceeds inflation: $s_t > \pi_t$,
    i.e.\
    \[
      t + 5 > t^2 - 25t + 155 \implies t^2 - 26t + 150 < 0.
    \]
    The roots of $t^2 - 26t + 150 = 0$ are $t = 13 \pm \sqrt{19}$, giving
    approximately $t \approx 8.64$ and $t \approx 17.36$. Real income rises
    for $t \in (8.64,\, 17.36)$ and falls outside this interval. Since real
    income is declining for $t < 8.64$ (including all of $t \leq 8$), it
    peaks at $t = 0$ --- which is not among options~A, B, or~C. \textbf{(D)
    None of the above} is correct.
  },
}
